%%%%%
%
%% Dans ce fichier, les �ventuels accents sont cod�s
%% en ISO-8859-15 (aka latin9)
%% La ligne suivante donne trois exemples 
%% e accent aigu : �, e accent grave : �, e-dans-l'o : �
\section{Transformation d'un profil de distance donn�es � partir d'une plateforme mobile dans un profil d'altitude }

\itemstyle{Op\'erateur~:} \operfmt{ProfilMobile2geo\_1D}
\index{ProfilMobile2geo\_1D}

%%%%%OBJET%%%%
\itemstyle{Objet~:}
Cet op�rateur transforme un profil de hauteurs donn�es � partir d'une plateforme mobile en un profil d'altitude dans le rep�re terrestre. Il utilise les angles d'attitudes et l'altitude  de la plateforme.
Cet op�rateur peut typiquement �tre utilis� pour des mesures lidar r�alis�es � partir d'une plateforme mobile.

%%%%%ARGUMENTS%%%%
\bigskip
%% Dans ce fichier, les �ventuels accents sont cod�s
%% en ISO-8859-15 (aka latin9)
%% La ligne suivante donne trois exemples 
%% e accent aigu : �, e accent grave : �, e-dans-l'o : �
%%%%%ENTREES%%%%
\begin{tabarg}
\multicolumn{3}{l}{\titretab {Entr\'ees}}\\
\hline
$X$ & Vecteur, \texttt{float[$N_{level}$]} & Signal mesur� \\
$D$ & Coeff., \texttt{int} & 1 si mesure vers le haut, -1 si mesure vers le bas  \\
$\phi$ & Vecteur, \texttt{float} & Roulis (degr�s)\\
$\theta$ & Vecteur, \texttt{float} & Tangage (degr�s)\\
$Alt$ & Vecteur, \texttt{float} & Altitude de la plateforme (m)\\
\end{tabarg}

\medskip
%%%%%SORTIES%%%%
\begin{tabarg}
\multicolumn{3}{l}{\titretab {Sorties}}\\
\hline
$Alt_X$ & Vecteur, \texttt{float[$N_{level}$]} & Altitude des mesures X (m)\\
\end{tabarg}



%%%%%FORMULE%%%%
\itemstyle{Formule}

%% Formule de l'operateur : ProfilMobile2geo_1D
%% Dans ce fichier, les �ventuels accents sont cod�s
%% en ISO-8859-15 (aka latin9)
%% La ligne suivante donne trois exemples 
%% e accent aigu : �, e accent grave : �, e-dans-l'o : �

Pour tous les niveaux de mesure, et pour chaque �chantillon 
temporel, l'altitude est calcul�e par formule suivante :

$$
 Alt_X = Alt + D * \cos(\theta)*\cos(\phi)* Z_X
$$



%%%%%AUTEUR%%%%
\itemstyle{Auteur~:}
CNRM/GMEI/TRAMM

%%%%%REFERENCE%%%%
\itemstyle{R\'ef\'erences~:}

